\documentclass[a4paper,12pt]{article}

% Pacchetti comuni
\usepackage[utf8]{inputenc}    % Codifica UTF-8
\usepackage{amsmath}           % Simboli matematici
\usepackage{amsfonts}          % Font matematici
\usepackage{amssymb}           % Simboli matematici
\usepackage{graphicx}          % Gestione delle immagini
\usepackage{hyperref}          % Collegamenti ipertestuali
\usepackage[a4paper,margin=1.2cm]{geometry}  % Margini stretti

\title{Formulario di Algebra Lineare}
\author{Nicolò Giuliani}
\date{}


\begin{document}
%\maketitle
\section*{Misc.}
\textbf{Vettori linearmente indipendenti}:
\begin{align*}
    &\alpha \vec{v}_1 + \beta \vec{v}_2 + \gamma \vec{v}_3 = \vec{0} \iff \alpha = \beta = \gamma = 0 \\
    &\iff \det([\vec{v}_1 \, \vec{v}_2 \, \vec{v}_3]) \neq 0  \\
    &\iff A = [\vec{v}_1 \, \vec{v}_2 \, \vec{v}_3] \xrightarrow{\text{Gauss}} A' \text{ (forma a scala)}
\end{align*}
\textbf{prodotto di matrici}
\begin{itemize}
    \item \( A_{m \times n} \cdot B_{n \times p} = C_{m \times p} \)
    \item Ogni elemento di \( C \) è dato da:
    \[
    c_{ij} = \sum_{k=1}^n a_{ik} b_{kj}
    \]
\end{itemize}
\section*{Proprietà Matrici (Sintesi)}

\begin{itemize}
  \item \(A + B = B + A\) (addizione commutativa)
  \item \(A(BC) = (AB)C\) (moltiplicazione associativa)
  \item \(A(B + C) = AB + AC\) (distributiva)
  \item \(AB \neq BA\) in generale (non commutativa)
  \item \(I\) matrice identità: \(AI = IA = A\)
  \item \((A^T)^T = A\), \((AB)^T = B^T A^T\)
  \item Se invertibile: \(AA^{-1} = I\), \((AB)^{-1} = B^{-1} A^{-1}\)
  \item \(\det(AB) = \det(A) \det(B)\), \(\det(A^T) = \det(A)\)
  \item Matrice ortogonale: \(Q^T Q = I\), \(Q^{-1} = Q^T\), \(\det(Q) = \pm 1\)
  \item Se \(Q,R\) ortogonali, allora \(QR\) ortogonale
\end{itemize}

\section*{Esercizio 1}
\begin{align*}
    F_k : \mathbb{R}^n \to \mathbb{R}^m\\
    F_k(x_1,x_2,\dots,x_n)=\begin{cases} 
        riga_1 \\
        riga_2 \\
        \vdots \\
        riga_m
    \end{cases}
\end{align*}
Sia $A \in M_{m\times n} (\mathbb{R})$: \\
\begin{align*}
    F(\vec{v})=A\times [\vec{v}]
\end{align*}
\subsection*{a)}
\textbf{Iniettiva}: $Ker(F_k)=\{0\} \iff det(A_k)\neq 0 $ \\
\textbf{Suriettiva}: $dim(Im)=dim(\mathbb{R}^{m})\iff det(A_k)\neq 0$ \\
\textbf{Isomorfismo/Invertibile}: $\exists F^{-1}_k \iff det(A_k)\neq0$ \\
\textbf{Teorema della dimensione}: $dim(Ker)=dim(\mathbb{R}^m)-dim(Im), dim(Ker)=\text{num. colonne} -rr(A')$ \\
\textbf{base di Ker}: $A'(\vec{x}) \to \begin{cases}
    x_1= \dots \\
    \vdots
\end{cases}$
\subsection*{b)}
\textbf{b.1.1) cerco i $k$ t.c. $\vec{v} \in Ker(F_k$)}: 
\begin{align*}
    \vec{v}=\{v_1,v_2,\dots,v_m\} \quad 
    A\vec{v}=0 \to \begin{cases}
        v_1=\dots \\
        v_2=\dots \\
        \vdots
    \end{cases}
    \to \boxed{k=\dots}
\end{align*}
\textbf{b.1.2) controllo se $\vec{v}\in Ker$}: 
\begin{align*}
    \vec{v}=\{v_1,v_2,\dots,v_m\}\in Ker(F)\iff Ker(F(\vec{v}))=\{0\}\\
    A\vec{v}=0 \to 
    \begin{cases}
        v_1(\dots)=0\\
        v_2(\dots)=0\\
        \vdots \\
        v_m(\dots)=0
    \end{cases}
    \to \begin{cases}
        v_1=0\\
        v_2=0\\
        \vdots \\
        v_m=0
    \end{cases} \to  (\top)\, \vec{v}\in Ker(F)
\end{align*}
\textbf{b.2) calcolo base di $Ker(F)$}:
\begin{align*}
    A\vec{x}=0 \to \begin{pmatrix}
        a_{11} & a_{12}\\
        a_{21} & a_{22}
    \end{pmatrix} 
    \to \begin{cases}
        a_{11}x_1 + a_{12}x_2 =0 \\
        \vdots
    \end{cases} \to \begin{cases}
        x_1=\dots \\
        x_2=\dots
    \end{cases} \\
    \to \beta_{Ker(F)}=\{(x_1,x_2,\dots)\}
\end{align*}
\textbf{b) 2.1 Controllo se \( \vec{v} \in \text{Im}(A) \)}:

\textbf{Metodo 1: Risolvere il sistema \( A \mathbf{x} = \vec{v} \)}

\[
A \mathbf{x} = \vec{v} \quad \text{dove} \quad A \in \mathbb{R}^{m \times n}, \quad \vec{v} \in \mathbb{R}^m
\]
\[
\to \begin{cases}
A_{1} \cdot x_1 + A_{2} \cdot x_2 + \dots + A_n \cdot x_n = v_1 \\
A_{1} \cdot x_1 + A_{2} \cdot x_2 + \dots + A_n \cdot x_n = v_2 \\
\vdots \\
A_{1} \cdot x_1 + A_{2} \cdot x_2 + \dots + A_n \cdot x_n = v_m
\end{cases}
\]
Se il sistema è \textbf{compatibile} (non hanno contraddizioni), allora \( \vec{v} \in \text{Im}(A) \). Se il sistema è \textbf{incompatibile}, allora \( \vec{v} \notin \text{Im}(A) \). \\
\textbf{Metodo 2: Usare il rango}: $rr(A)=rr(A\mid \vec{v}) \iff \vec{v} \in Im(A)$  \\
\vspace{1em}
\textbf{b.3) Calcolare le coordinate di $\vec{v}$ rispetto a $\beta$:}
\begin{align*}
    &\beta = \{\vec{b}_1, \vec{b}_2\} \\
    &[\vec{v}]_{\beta} = \begin{bmatrix} y_1 \\ y_2 \end{bmatrix} \\
    &\Rightarrow \vec{v} = y_1 \vec{b}_1 + y_2 \vec{b}_2\\
    \to \begin{pmatrix}
        v_1\\
        v_2
    \end{pmatrix}=y_1\begin{pmatrix}
        [\vec{b_1}]\\
        \vdots
    \end{pmatrix}
    +y_2\begin{pmatrix}
        [\vec{b_2}]\\
        \vdots
    \end{pmatrix}
\end{align*}
\subsection*{c)}
\textbf{c.1.1) Scrivere le equazioni cartesiane di $\ker(F_k)$}:

\begin{align*}
    &A\vec{x} = \vec{0} \Rightarrow 
    \begin{cases}
        a_{11}x_1 + a_{12}x_2 + \dots + a_{1n}x_n = 0 \\
        a_{21}x_1 + a_{22}x_2 + \dots + a_{2n}x_n = 0 \\
        \vdots \\
        a_{m1}x_1 + a_{m2}x_2 + \dots + a_{mn}x_n = 0
    \end{cases} \\
    &\text{Semplificando il sistema otteniamo ad esempio:} \\
    &x_i = 2x_k + 4x_p \iff \boxed{x_i - 2x_k - 4x_p = 0}
\end{align*}

\textbf{d.1.2) Scrivere le equazioni cartesiane di $\text{Im}(F_{(k)})$, (al variare di $k$):}
\begin{align*}
    &A_k\vec{x} = \vec{y} \Rightarrow 
    \begin{cases}
        a_{11}(k)x_1 + a_{12}(k)x_2 + \dots + a_{1n}(k)x_n = y_1 \\
        a_{21}(k)x_1 + a_{22}(k)x_2 + \dots + a_{2n}(k)x_n = y_2 \\
        \vdots \\
        a_{m1}(k)x_1 + a_{m2}(k)x_2 + \dots + a_{mn}(k)x_n = y_m
    \end{cases} \\
    &\boxed{y_1 = a_{11}(k)x_1 + a_{12}(k)x_2 + \dots + a_{1n}(k)x_n, \quad y_2 = a_{21}(k)x_1 + a_{22}(k)x_2 + \dots + a_{2n}(k)x_n, \dots}
\end{align*}
probabilmente meglio:
\begin{align}
    A^t \to A^{t'} \\
    \text{Semplificando, tenendo gli $x_i=\top$  e sostituendo gli $x_j=0$:} \\ 
    \to \begin{cases}
        \boxed{x_1-x_k=0} \\
        x_j=0 \\
        x_k=\top
    \end{cases}
\end{align}


\pagebreak
\section*{Esercizio 2}
\begin{align*}
    T_k : \mathbb{R}^n \to \mathbb{R}^n
\end{align*}
Sia \( A \) la matrice associata a \( T_k \) rispetto alla base canonica (in dominio e codominio).
\begin{align*}
    A_k = \begin{pmatrix}
        [T_k(e_1)] & [T_k(e_2)] & \dots & [T_k(e_n)] \\
        \vdots & \vdots & & \vdots
    \end{pmatrix}
\end{align*}

\textbf{a.1.1) Valori di \( k \) tali che \( T_k \) è diagonalizzabile:}

\begin{align*}
    T_k \text{ è diagonalizzabile} \iff \exists \beta = \{ \vec{a_1}, \vec{a_2}, \dots, \vec{a_n} \} \quad \text{con} \quad \vec{a_i} = \text{span}\left( [A] - \lambda_i I = 0 \right)
\end{align*}

I valori di \( k \) comportano il controllo delle seguenti condizioni:

- \textbf{MA} (molteplicità algebrica): il grado del polinomio caratteristico associato agli autovalori \( \lambda_i \). \\
- \textbf{MG} (molteplicità geometrica): la dimensione del nucleo \( \text{Ker}(A - \lambda_i I) \), che corrisponde al numero di autovettori indipendenti associati a ciascun autovalore. \\

Per determinare quando \( T_k \) è diagonalizzabile, bisogna verificare che la molteplicità geometrica sia uguale alla molteplicità algebrica per ogni autovalore \( \lambda_i \).

\textbf{Esempio di tabella per analizzare i valori di \( k \):}

\begin{tabular}{|c|c|c|c|}
    \hline
    \( \lambda_i \) & MA (Molteplicità Algebrica) & MG (Molteplicità Geometrica) & \( k \) \\
    \hline
    0 & 2 & 2 & \( k=0 \) \\
    \hline
    1 & 3 & 3 & \( k \neq 0 \) \\
    \hline
    2 & 1 & 1 & \( k=2 \) \\
    \hline
\end{tabular}
\vspace{1em}

\textbf{b.1.1) Valori di $k$ t.c. $\vec{v}$ e' autovettore di $F_k$}:
\begin{align}
    F_k(\vec{v})=\lambda[\vec{v}] \\
    \begin{cases}
        riga_1=\lambda v_1 \\
        riga_2=\lambda v_2\\
        \vdots \\
        riga_n=\lambda v_n
    \end{cases}
\end{align}
\textbf{c.1.1) Matrice associata nel dominio $(C)$ e codominio ($\beta$)}: 
\begin{align}
    [M]_{\beta}^{C} &= P_{\beta}^{-1} \cdot A \cdot P_C \\
    P_C &= I_n \quad \text{(base canonica)} \Rightarrow A \cdot P_C = A \\
    \beta &= \{ \vec{b}_1,\ \vec{b}_2,\ \vec{b}_3 \} \Rightarrow P_{\beta} = [\vec{b}_1\ \vec{b}_2\ \vec{b}_3]
\end{align}


\pagebreak
\section*{Esercizio 3}
Due vettori (o più) sono \textbf{a due a due ortogonali} se per ogni coppia distinta \(i \neq j\) vale:
\[
\mathbf{v}_i \cdot \mathbf{v}_j = 0
\]
dove \(\mathbf{v}_i, \mathbf{v}_j\) sono vettori nello spazio considerato. \\
\textbf{a.0)}
\begin{align*}
    A_{ortogonale}\in \mathbb{M}_n(\mathbb{R}) \to A^T A = I_n \iff A^T=A^{-1}
\end{align*}
$det(A)=\pm 1$ \\
In $\mathbb{R}^3$:
\[
U = \langle \vec{u}_1, \vec{u}_2 \rangle
\]

\textbf{a.1.0) Costruzione di un vettore $\vec{v}$ ortogonale a $\langle n, g \rangle$:}

\begin{align*}
\mathrm{proj}_{\langle n,g \rangle}(\vec{v}) 
&= \frac{\langle \vec{v},n \rangle}{\langle n,n \rangle}n 
+ \frac{\langle \vec{v},g \rangle}{\langle g,g \rangle}g \\
&= \begin{pmatrix}
\dots & \dots & \dots
\end{pmatrix}
\end{align*}

\textbf{a.1.1) Costruzione del vettore ortogonale $\vec{v}_\perp$:}

\begin{align*}
\vec{v}_\perp 
&= \vec{v} - \mathrm{proj}_{\langle n, g \rangle}(\vec{v}) \\
&= \begin{pmatrix}
\cdots \\ \cdots \\ \cdots
\end{pmatrix}
\end{align*}

\textbf{a.1.2) Calcolo di una base ortogonale di $U$:}

\begin{align*}
\beta = \{\vec{v}_1, \vec{v}_2\} = \left\{\vec{v}_1, \vec{u}_2 - \mathrm{proj}_{\vec{v}_1}(\vec{u}_2) \right\}
\end{align*}

\textbf{a.1.3) Calcolo di una base di $U^\perp$:}

\begin{align*}
\begin{cases}
\langle \vec{v}, \vec{v}_1 \rangle = 0 \\
\langle \vec{v}, \vec{v}_2 \rangle = 0
\end{cases}
\Rightarrow
\vec{v} = \begin{pmatrix}
x \\ y \\ z
\end{pmatrix},\quad 
\vec{v}_1 = \begin{pmatrix}
[u_1] \\ \vdots
\end{pmatrix},\quad 
\vec{v}_2 = \begin{pmatrix}
[u_2] \\ \vdots
\end{pmatrix} \\
\Rightarrow
\begin{cases}
x - y = 0 \\
x - z = 0
\end{cases}
\Rightarrow
x = y = z \Rightarrow 
\vec{v} = x \begin{pmatrix}
1 \\ 1 \\ 1
\end{pmatrix}
\Rightarrow 
\tilde{\beta} = \left\{ \begin{pmatrix}
1 \\ 1 \\ 1
\end{pmatrix} \right\}
\end{align*}

\textit{Nota:} $\dim(\mathbb{R}^3) = \dim(\beta) + \dim(\tilde{\beta})$

\textbf{b.1.1) Applicazione lineare della proiezione ortogonale su $\langle n \rangle$:}

\begin{align*}
F : \mathbb{R}^3 \to \mathbb{R}^3 \\
A = \begin{pmatrix}
\left[ \mathrm{proj}_{\langle n \rangle}(e_1) \right] & 
\left[ \mathrm{proj}_{\langle n \rangle}(e_2) \right] & 
\left[ \mathrm{proj}_{\langle n \rangle}(e_3) \right]
\end{pmatrix}
\end{align*}

\textbf{c.1.1) Equazioni cartesiane di un sottospazio $W \subseteq \mathbb{R}^4$:}

\begin{align*}
W = \langle \vec{w}_1, \vec{w}_2 \rangle 
= \lambda \vec{w}_1 + \gamma \vec{w}_2 
= (\lambda w_{11} + \gamma w_{12}, \lambda w_{21} + \gamma w_{22}, \lambda w_{31} + \gamma w_{32}, \lambda w_{41} + \gamma w_{42}) \\
\Rightarrow A \mid \vec{x} \Rightarrow \begin{cases}
\lambda w_{11} + \gamma w_{12} = x_1 \\
\lambda w_{21} + \gamma w_{22} = x_2 \\
\lambda w_{31} + \gamma w_{32} = x_3 \\
\lambda w_{41} + \gamma w_{42} = x_4
\end{cases}
\end{align*}
\textbf{c.1.2) controllo $\vec{v}\in W\subseteq \mathbb{R}^n$}
\begin{align*}
    A\lambda_i=\vec{v}
\end{align*}
se ci sono soluzioni $\vec{v} \in W$.
\pagebreak
\section*{Esercizio 4 — Congruenze lineari}

Una \textbf{congruenza lineare} è un'equazione del tipo:
\[
ax \equiv b \pmod{n}
\]
Essa ammette soluzioni intere se e solo se:
\[
\gcd(a, n) \mid b,
\]
cioè il massimo comun divisore tra \(a\) e \(n\) deve dividere \(b\).

\subsection*{Esempio svolto}

Risolvere la seguente congruenza lineare:
\[
38x \equiv 5 \pmod{87}
\]

\paragraph{1. Calcolo del massimo comun divisore}
Applichiamo l'algoritmo di Euclide:
\[
\begin{aligned}
87 &= 2 \cdot 38 + 11 \\
38 &= 3 \cdot 11 + 5 \\
11 &= 2 \cdot 5 + 1 \\
5 &= 5 \cdot \boxed{1} + 0
\end{aligned}
\]
Dunque:
\[
\gcd(38, 87) = 1
\]
Poiché \(1 \mid 5\), la congruenza ammette soluzioni intere.

\paragraph{2. Inverso moltiplicativo di 38 modulo 87}
Applichiamo l'algoritmo di Euclide esteso:
\[
\begin{aligned}
1 &= 11 - 2 \cdot 5 \\
  &= 11 - 2(38 - 3 \cdot 11) = 7 \cdot 11 - 2 \cdot 38 \\
  &= 7(87 - 2 \cdot 38) - 2 \cdot 38 = 7 \cdot 87 - 16 \cdot 38 \\
&\boxed{1 = 7 \cdot 87 - 16 \cdot 38}
\end{aligned}
\]
Dunque:
\[
1 = 7 \cdot 87 - 16 \cdot 38 \Rightarrow -16 \cdot 38 \equiv 1 \pmod{87}
\]
Quindi l'inverso di \(38\) modulo \(87\) è \(-16\): 
\begin{align*}
    \boxed{38^{-1} = -16 \text{ (mod 87)}}
\end{align*}

\paragraph{3. Soluzione generale}
Moltiplichiamo entrambi i membri della congruenza per l'inverso:
\begin{align*}
    38x &\equiv 5 \pmod{87} \\
    x &\equiv -16 \cdot 5 \pmod{87} \\
    x &\equiv (87-16) \cdot 5 \pmod{87} \\
    x &\equiv 71 \cdot 5 \pmod{87} \\
    x &\equiv 355 \pmod{87}
\end{align*}
Riduciamo il risultato modulo \(87\):
\[
355 \equiv \boxed{7} \pmod{87} \iff 355-87*4=\boxed{7}
\]

Pertanto, la soluzione generale è:
\[
\boxed{x = 7 + 87k, \quad k \in \mathbb{Z}}
\]



\section*{2. Classi di resto}

La \textbf{classe di resto} di un numero intero \(a\) modulo \(n\), indicata con \([a]_n\), è l’insieme di tutti i numeri interi che sono congrui ad \(a\) modulo \(n\). In simboli:
\[
[a]_n = \{x \in \mathbb{Z} \mid x \equiv a \pmod{n} \}
\]
Due numeri \(a\) e \(b\) appartengono alla stessa classe modulo \(n\) se:
\[
n \mid (a - b) \quad \Longleftrightarrow \quad a \equiv b \pmod{n}
\]

\subsection*{Esempio}
Mostriamo che \(17\) e \(5\) appartengono alla stessa classe modulo \(12\).

Osserviamo che:
\[
17 - 5 = 12 \Rightarrow 12 \mid (17 - 5)
\]
Quindi:
\[
17 \equiv 5 \pmod{12}
\]
Dunque:
\[
[17]_{12} = [5]_{12}
\]

\section*{3. Inclusione tra classi di resto}

Dire che:
\[
[a]_m \subseteq [b]_n
\]
significa che ogni numero intero congruo ad \(a\) modulo \(m\) è anche congruo a \(b\) modulo \(n\). In altre parole:
\[
\forall x \in \mathbb{Z},\quad x \equiv a \pmod{m} \Rightarrow x \equiv b \pmod{n}
\]

Per verificare ciò, si risolve il seguente sistema di congruenze:
\[
\begin{cases}
x \equiv a \pmod{m} \\
x \equiv b \pmod{n}
\end{cases}
\]
e si verifica se \emph{ogni soluzione} di \(x \equiv a \pmod{m}\) è anche soluzione di \(x \equiv b \pmod{n}\).

Se il sistema ha: \\
- più soluzioni diverse $\to$ allora \([a]_m \not\subseteq [b]_n\); \\
- una soluzione sola $\to$ allora bisogna verificare se tutte le soluzioni di \(x \equiv a \pmod{m}\) soddisfano anche \(x \equiv b \pmod{n}\); \\
- nessuna soluzione $\to$ allora \([a]_m \not\subseteq [b]_n\). \\

\subsection*{Esempio}
Verificare se:
\[
[-7]_{25} \subseteq [18]_{50}
\]

\paragraph{1. Scriviamo la congruenza}  
Dalla congruenza \(x \equiv -7 \pmod{25}\), otteniamo la forma esplicita:
\[
x = 25k - 7, \quad k \in \mathbb{Z}
\]

\paragraph{2. Sostituiamo nella seconda congruenza}  
Sostituiamo \(x = 25k - 7\) nella congruenza \(x \equiv 18 \pmod{50}\):
\[
25k - 7 \equiv 18 \pmod{50}
\]

\paragraph{3. Risolviamo la congruenza}
Portiamo i termini noti a destra:
\[
25k \equiv 25 \pmod{50}
\Rightarrow 25k - 25 \equiv 0 \pmod{50}
\Rightarrow 25(k - 1) \equiv 0 \pmod{50}
\]

Poiché \(25 \mid 50\), possiamo dividere ambo i membri per \(25\):
\[
k - 1 \equiv 0 \pmod{2} \Rightarrow k \equiv 1 \pmod{2}
\]

\paragraph{4. Conclusione}
Abbiamo trovato che solo i \(k\) dispari (ovvero \(k \equiv 1 \pmod{2}\)) soddisfano la seconda congruenza. Quindi:
\[
x = 25k - 7 \equiv 18 \pmod{50} \quad \text{solo per alcuni } k
\]
Non tutte le soluzioni della prima congruenza sono anche soluzioni della seconda.

\[
\boxed{[-7]_{25} \not\subseteq [18]_{50}}
\]


\end{document}